\clearpage
\begin{trapbox}
\textbf{\textcolor{CTrap}{最容易错的八个地方}}
\textbf{\textcolor{CTrap}{易错1：“有一个周期”被偷换成“最小正周期”。}}
\begin{center}\begin{tikzpicture}[font=\small,>=stealth,line join=round]
\draw[gray!40,rounded corners=3pt] (0,0) rectangle (14.2,8.3);\node at (7.1,7.700000000000001) {例2：用周期4只能列出偶数轴，会漏掉奇数轴};
\begin{scope}[shift={(2.1714285714285717,3.8222222222222224)},x=1.0714285714285714cm,y=1.277777777777778cm]\draw[->,gray] (-1,0)--(4.6,0) node[right]{$x$};\draw[->,gray] (0,-1.7)--(0,1.9) node[above]{$y$};\draw[blue,thick,solid,domain=-0.7:4.2,samples=180] plot(\x,{cos(deg(pi*\x))});\foreach \a in {0,2,4}{\draw[violet,dashed,thick] (\a,-1.3)--(\a,1.3);}\foreach \a in {1,3}{\draw[orange!85!black,dotted,thick] (\a,-1.3)--(\a,1.3);}\end{scope}\node[anchor=west,align=left,text width=5.0cm] at (8.4,4.0) {$f=\cos\pi x$\\[9pt]已知一个周期 $T=4$\\[9pt]紫线：公式保证 $x=2k$\\[9pt]橙点线：$x=1,3$ 也是真轴\\[9pt]最小周期2才列齐 $x=k$};
\node[text width=12.6cm,align=center] at (7.1,.65) {两类竖线都是真对称轴；已知的一个周期越大，用它列出的一组位置可能越稀。};
\end{tikzpicture}\end{center}



例如 $\cos(\pi x)$ 的周期可以取 $4$，但最小正周期是 $2$。从 $T=4$ 只能推出 $x=2k$ 这组对称轴，不能说它包含全部对称轴。



\end{trapbox}
\clearpage
\begin{trapbox}
\textbf{\textcolor{CTrap}{易错2：把“半周期出现对称轴”理解成“半周期仍重复”。}}
\begin{center}\begin{tikzpicture}[font=\small,>=stealth,line join=round]
\draw[gray!40,rounded corners=3pt] (0,0) rectangle (14.2,8.3);\node at (7.1,7.700000000000001) {周期2的余弦：轴每隔1出现，中心也每隔1出现};
\begin{scope}[shift={(2.3,3.8222222222222224)},x=1.2cm,y=1.277777777777778cm]\draw[->,gray] (-1,0)--(4,0) node[right]{$x$};\draw[->,gray] (0,-1.7)--(0,1.9) node[above]{$y$};\draw[blue,thick,solid,domain=-0.7:3.7,samples=180] plot(\x,{cos(deg(pi*\x))});\fill[orange!85!black] (0.5,0) circle(2.2pt);\fill[orange!85!black] (1.5,0) circle(2.2pt);\fill[orange!85!black] (2.5,0) circle(2.2pt);\fill[orange!85!black] (3.5,0) circle(2.2pt);\foreach \a in {0,1,2,3}{\draw[violet,dashed] (\a,-1.3)--(\a,1.3);\node[above] at (\a,1.35){$\a$};}\end{scope}\node[anchor=west,align=left,text width=5.0cm] at (8.4,4.0) {$f=\cos\pi x,\ T_0=2$\\[9pt]紫线：轴 $x=k$\\[9pt]橙点：中心 $(\frac12+k,0)$\\[9pt]同类间隔 $\frac{T_0}2=1$\\[9pt]但 $f(0)=1\ne f(1)=-1$};
\node[text width=12.6cm,align=center] at (7.1,.65) {半周期是同类对称位置的间隔，不表示每个函数值都在半周期后重复。};
\end{tikzpicture}\end{center}



对于 $\cos(\pi x)$，$x=0$ 和 $x=1$ 都是对称轴，可是 $f(0)=1$，$f(1)=-1$。\textbf{半周期不是周期，不能随意写 $f(x+T/2)=f(x)$。}



\end{trapbox}
\clearpage
\begin{trapbox}
\textbf{\textcolor{CTrap}{易错3：求对称中心时把纵坐标也加了 $T/2$。}}
\begin{center}\begin{tikzpicture}[font=\small,>=stealth,line join=round]
\draw[gray!40,rounded corners=3pt] (0,0) rectangle (14.2,8.3);\node at (7.1,7.700000000000001) {例3：中心横向移3，高度始终保留5};
\begin{scope}[shift={(1.4000000000000001,1.8921052631578945)},x=0.6cm,y=0.4842105263157894cm]\draw[->,gray] (-0.5,0)--(9.5,0) node[right]{$x$};\draw[->,gray] (0,-0.5)--(0,9) node[above]{$y$};\draw[teal,thick,solid,domain=0:9,samples=180] plot(\x,{5+2*sin(deg(pi*(\x-2)/3))});\fill[orange!85!black] (2,5) circle(2.2pt);\fill[orange!85!black] (5,5) circle(2.2pt);\fill[orange!85!black] (8,5) circle(2.2pt);\fill[orange!85!black] (3,6.732050807568877) circle(2.2pt);\fill[orange!85!black] (7,3.267949192431123) circle(2.2pt);\draw[gray,dashed] (0,5)--(9,5);\foreach \v in {2,5,8}{\node[below=4pt] at (\v,0){$\v$};}\end{scope}\node[anchor=west,align=left,text width=5.0cm] at (8.4,4.0) {$f=5+2\sin\frac{\pi(x-2)}3$\\[9pt]$T=6$，中心 $(2,5),(5,5)$\\[9pt]$f(5)=f(8)=5$\\[9pt]$f(3)+f(7)=10$\\[9pt]纵坐标不增加 $T/2$};
\node[text width=12.6cm,align=center] at (7.1,.65) {橙点包括中心与绕中心配对的点；相邻中心横坐标相差3，高度都为5。};
\end{tikzpicture}\end{center}



正确形式是 $(b+kT/2,c)$，而不是 $(b+kT/2,c+kT/2)$。一个周期只对应\textbf{水平方向}的重复。



\end{trapbox}
\clearpage
\begin{trapbox}
\textbf{\textcolor{CTrap}{易错4：两条对称轴推出的周期少乘了 $2$。}}
\begin{center}\begin{tikzpicture}[font=\small,>=stealth,line join=round]
\draw[gray!40,rounded corners=3pt] (0,0) rectangle (14.2,8.3);\node at (7.1,7.700000000000001) {证明一：右点先移回一个周期，再沿旧轴配对};
\begin{scope}[shift={(2.3857142857142857,3.8222222222222224)},x=0.8571428571428571cm,y=1.277777777777778cm]\draw[->,gray] (-1.5,0)--(5.5,0) node[right]{$x$};\draw[->,gray] (0,-1.7)--(0,1.9) node[above]{$y$};\draw[blue,thick,solid,domain=-1:5,samples=180] plot(\x,{cos(deg(pi*(\x-1)/2))});\fill[orange!85!black] (-0.5,-0.7071067811865476) circle(2.2pt);\fill[orange!85!black] (2.5,-0.7071067811865476) circle(2.2pt);\fill[orange!85!black] (3.5,-0.7071067811865476) circle(2.2pt);\foreach \a in {1,3}{\draw[violet,dashed] (\a,-1.25)--(\a,1.25);\node[above] at (\a,1.3){$x=\a$};}\foreach \X/\L in {-.5/{R-T},2.5/L,3.5/R}{\draw[gray,dashed] (\X,-.70710678)--(\X,-1.2);\node[below=4pt] at (\X,-1.4){$\L$};}\end{scope}\node[anchor=west,align=left,text width=5.0cm] at (8.4,4.0) {$T=4,\ a=1,\ t=\frac12$\\[9pt]新轴 $a+T/2=3$\\[9pt]$R=3.5,\ R-T=-0.5$\\[9pt]$L=2.5$\\[9pt]三处都为 $-\frac{\sqrt2}2$};
\node[text width=12.6cm,align=center] at (7.1,.65) {图取一组具体距离；证明中t任意，平移一次与沿旧轴配对就得到新轴。};
\end{tikzpicture}\end{center}



如果两轴在 $x=1$ 与 $x=3$，间距是 $2$，由两次对折得到一个周期是 $4$，不是直接把 $2$ 当周期。



\end{trapbox}
\clearpage
\begin{trapbox}
\textbf{\textcolor{CTrap}{易错5：把求出的 $2|a_1-a_2|$ 说成最小正周期。}}
\begin{center}\begin{tikzpicture}[font=\small,>=stealth,line join=round]
\draw[gray!40,rounded corners=3pt] (0,0) rectangle (14.2,8.3);\node at (7.1,7.700000000000001) {例2：用周期4只能列出偶数轴，会漏掉奇数轴};
\begin{scope}[shift={(2.1714285714285717,3.8222222222222224)},x=1.0714285714285714cm,y=1.277777777777778cm]\draw[->,gray] (-1,0)--(4.6,0) node[right]{$x$};\draw[->,gray] (0,-1.7)--(0,1.9) node[above]{$y$};\draw[blue,thick,solid,domain=-0.7:4.2,samples=180] plot(\x,{cos(deg(pi*\x))});\foreach \a in {0,2,4}{\draw[violet,dashed,thick] (\a,-1.3)--(\a,1.3);}\foreach \a in {1,3}{\draw[orange!85!black,dotted,thick] (\a,-1.3)--(\a,1.3);}\end{scope}\node[anchor=west,align=left,text width=5.0cm] at (8.4,4.0) {$f=\cos\pi x$\\[9pt]已知一个周期 $T=4$\\[9pt]紫线：公式保证 $x=2k$\\[9pt]橙点线：$x=1,3$ 也是真轴\\[9pt]最小周期2才列齐 $x=k$};
\node[text width=12.6cm,align=center] at (7.1,.65) {两类竖线都是真对称轴；已知的一个周期越大，用它列出的一组位置可能越稀。};
\end{tikzpicture}\end{center}



例如 $\cos(\pi x)$ 的 $x=0$ 与 $x=2$ 都是轴，两倍距离得到 $4$，但真正最小正周期是 $2$。



\end{trapbox}
\clearpage
\begin{trapbox}
\textbf{\textcolor{CTrap}{易错6：把一个对称中心当成对称轴来用。}}
\begin{center}\begin{tikzpicture}[font=\small,>=stealth,line join=round]
\draw[gray!40,rounded corners=3pt] (0,0) rectangle (14.2,8.8);\node at (7.1,8.200000000000001) {中心两边的值相加为6，轴两边的值彼此相等};
\node at (3.1500000000000004,7.3) {$f=3+\sin\pi(x-1)$};\begin{scope}[shift={(1.925,2.8454545454545457)},x=1.225cm,y=0.6909090909090909cm]\draw[->,gray] (-1,0)--(3,0) node[right]{$x$};\draw[->,gray] (0,-0.5)--(0,5) node[above]{$y$};\draw[teal,thick,solid,domain=-0.5:2.5,samples=180] plot(\x,{3+sin(deg(pi*(\x-1)))});\fill[orange!85!black] (0.5,2) circle(2pt);\fill[orange!85!black] (1.5,4) circle(2pt);\fill[orange!85!black] (1,3) circle(2pt);\draw[gray,dashed] (-.7,3)--(2.7,3);\end{scope}\node[align=center,text width=5.8cm] at (3.1500000000000004,1.7) {中心 $(1,3)$：两值为2和4\\相加6，但不相等};\node at (9.95,7.3) {$g=3+\cos\pi(x-1)$};\begin{scope}[shift={(8.725,2.8454545454545457)},x=1.225cm,y=0.6909090909090909cm]\draw[->,gray] (-1,0)--(3,0) node[right]{$x$};\draw[->,gray] (0,-0.5)--(0,5) node[above]{$y$};\draw[blue,thick,solid,domain=-0.5:2.5,samples=180] plot(\x,{3+cos(deg(pi*(\x-1)))});\fill[orange!85!black] (0.5,3) circle(2pt);\fill[orange!85!black] (1.5,3) circle(2pt);\draw[violet,dashed] (1,-.3)--(1,4.7);\end{scope}\node[align=center,text width=5.8cm] at (9.95,1.7) {轴 $x=1$：两值都为3\\左右等距点彼此同高};
\node[text width=12.6cm,align=center] at (7.1,.65) {两图同尺度；须按题设判断轴或中心，不能只凭两点等高猜对称类型。};
\end{tikzpicture}\end{center}



对称轴两边的函数值是\textbf{相等}：$f(a+t)=f(a-t)$；对称中心两边函数值是\textbf{相加为 $2c$}：$f(b+t)+f(b-t)=2c$。例如以 $(b,3)$ 为中心，函数值之和是 $6$，不是一定相等，也不是一定和为 $0$。



\end{trapbox}
\clearpage
\begin{trapbox}
\textbf{\textcolor{CTrap}{易错7：未经说明，断言对称中心的“全部”位置。}}
\begin{center}\begin{tikzpicture}[font=\small,>=stealth,line join=round]
\draw[gray!40,rounded corners=3pt] (0,0) rectangle (14.2,8.8);\node at (7.1,8.200000000000001) {全部中心的补充证明：有固定上下界，才排得掉高度漂移};
\node at (3.1500000000000004,7.3) {有界：$f=\cos\pi x$};\begin{scope}[shift={(1.5166666666666666,3.45)},x=0.8166666666666668cm,y=0.475cm]\draw[->,gray] (-1,0)--(5,0) node[right]{$x$};\draw[->,gray] (0,-2)--(0,6) node[above]{$y$};\draw[blue,thick,solid,domain=-0.7:5,samples=180] plot(\x,{cos(deg(pi*\x))});\fill[orange!85!black] (0.5,0) circle(2pt);\fill[orange!85!black] (1.5,0) circle(2pt);\fill[orange!85!black] (2.5,0) circle(2pt);\draw[gray,dashed] (-.8,1)--(4.8,1);\draw[gray,dashed] (-.8,-1)--(4.8,-1);\end{scope}\node[align=center,text width=5.8cm] at (3.1500000000000004,1.7) {所有值始终在 $[-1,1]$\\中心的高度不能不断增加};\node at (9.95,7.3) {漂移示例：$g=x$};\begin{scope}[shift={(8.316666666666666,3.45)},x=0.8166666666666668cm,y=0.475cm]\draw[->,gray] (-1,0)--(5,0) node[right]{$x$};\draw[->,gray] (0,-2)--(0,6) node[above]{$y$};\draw[orange!85!black,thick,solid,domain=-0.7:5,samples=180] plot(\x,{\x});\fill[orange!85!black] (0,0) circle(2pt);\fill[orange!85!black] (2,2) circle(2pt);\fill[orange!85!black] (1,1) circle(2pt);\fill[orange!85!black] (5,5) circle(2pt);\end{scope}\node[align=center,text width=5.8cm] at (9.95,1.7) {中心 $(0,0),(2,2)$\\迭代后高度从1变5再变9};
\node[text width=12.6cm,align=center] at (7.1,.65) {右图不是周期函数；有界性排除高度漂移。};
\end{tikzpicture}\end{center}



仅知道 $T$ 是某个周期、$(b,c)$ 是某个中心时，只能说 $(b+kT/2,c)$\textbf{肯定都是}中心。若另有“最小正周期”和“函数值有固定上下界”，才能证明这一组就是全部中心。不能把“全部轴”的证明无条件照搬过来。



\end{trapbox}
\clearpage
\begin{trapbox}
\textbf{\textcolor{CTrap}{易错8：把“四分之一周期”当作没有条件的公式。}}
\begin{center}\begin{tikzpicture}[font=\small,>=stealth,line join=round]
\draw[gray!40,rounded corners=3pt] (0,0) rectangle (14.2,8.3);\node at (7.1,7.700000000000001) {轴到中心不一定最近：也可能跨3个四分之一周期};
\begin{scope}[shift={(2.158823529411765,3.8222222222222224)},x=1.7647058823529411cm,y=1.277777777777778cm]\draw[->,gray] (-0.6,0)--(2.8,0) node[right]{$x$};\draw[->,gray] (0,-1.7)--(0,1.9) node[above]{$y$};\draw[blue,thick,solid,domain=-0.3:2.5,samples=180] plot(\x,{cos(deg(pi*\x))});\fill[orange!85!black] (0,1) circle(2.2pt);\fill[orange!85!black] (0.5,0) circle(2.2pt);\fill[orange!85!black] (1,-1) circle(2.2pt);\fill[orange!85!black] (1.5,0) circle(2.2pt);\draw[violet,dashed] (0,-1.3)--(0,1.3);\foreach \v/\L in {.5/\frac12,1/1,1.5/\frac32,2/2}{\node[below=4pt] at (\v,-1.35){$\L$};}\end{scope}\node[anchor=west,align=left,text width=5.0cm] at (8.4,4.0) {$T_0=2$，轴 $x=0$\\[9pt]中心 $(\frac12,0),(\frac32,0)$\\[9pt]$\frac12=T_0/4$\\[9pt]$\frac32=3T_0/4$\\[9pt]一般为奇数个 $T_0/4$};
\node[text width=12.6cm,align=center] at (7.1,.65) {四分之一周期关系使用最小正周期；不能把较远中心也误当最近的中心。};
\end{tikzpicture}\end{center}



$|b-a|=T_0/4$ 并不总成立：从一条轴到某个较远的中心，可能是 $3T_0/4,5T_0/4,\ldots$。正确说法是\textbf{横向有符号距离 $b-a$ 为奇数个 $T_0/4$}，还必须有\textbf{最小正周期 $T_0$} 与非恒定函数的前提。



\medskip
\end{trapbox}\clearpage\begin{trapbox}
\begin{center}\begin{tikzpicture}[font=\small,>=stealth,line join=round]
\draw[gray!40,rounded corners=3pt] (0,0) rectangle (14.2,8.8);\node at (7.1,8.200000000000001) {一对数值符合条件，不能代表整幅图像};
\node at (3.1500000000000004,7.3) {$u=x-3,\ f=u^3+u^2-u$};\begin{scope}[shift={(1.1083333333333334,3.26)},x=0.8166666666666668cm,y=0.2533333333333333cm]\draw[->,gray] (-0.5,0)--(5.5,0) node[right]{$x$};\draw[->,gray] (0,-3)--(0,12) node[above]{$y$};\draw[blue,thick,solid,domain=1:5,samples=180] plot(\x,{(\x-3)^3+(\x-3)^2-(\x-3)});\fill[orange!85!black] (2,1) circle(2pt);\fill[orange!85!black] (4,1) circle(2pt);\fill[orange!85!black] (2.5,0.625) circle(2pt);\fill[orange!85!black] (3.5,-0.125) circle(2pt);\draw[red,dotted] (3,-2.5)--(3,10.5);\end{scope}\node[align=center,text width=5.8cm] at (3.1500000000000004,1.7) {$f(2)=f(4)=1$\\但 $f(2.5)\ne f(3.5)$};\node at (9.95,7.3) {$g=(x-3)^2$};\begin{scope}[shift={(7.908333333333333,3.26)},x=0.8166666666666668cm,y=0.2533333333333333cm]\draw[->,gray] (-0.5,0)--(5.5,0) node[right]{$x$};\draw[->,gray] (0,-3)--(0,12) node[above]{$y$};\draw[teal,thick,solid,domain=0:5.5,samples=180] plot(\x,{(\x-3)^2});\fill[orange!85!black] (1,4) circle(2pt);\fill[orange!85!black] (5,4) circle(2pt);\fill[orange!85!black] (3,0) circle(2pt);\end{scope}\node[align=center,text width=5.8cm] at (9.95,1.7) {$g(1)+g(5)=8$\\但 $g(3)=0\ne4$};
\node[text width=12.6cm,align=center] at (7.1,.65) {左图不能认定轴x=3，右图不能认定中心(3,4)；等式必须对所有距离成立。};
\end{tikzpicture}\end{center}

\textbf{\textcolor{CTrap}{特别留心：公式是对任意实数成立，不是只看一两组点}}

例如检查 $f(2)=f(4)$，不能据此就说 $x=3$ 是整幅图像的对称轴；还要有
\[
f(3+t)=f(3-t)\quad\text{对所有 $t$ 成立}.
\]
同理，验证 $f(1)+f(5)=8$，也不能只凭这一对值就认定 $(3,4)$ 是图像的对称中心。

\medskip
\textbf{\textcolor{CTrap}{避免不必要的复杂性}}

整个结论只需会函数图像的对折、点的对称、周期与简单代数式变形。考试中如果已经知道一条轴和一个周期，\textbf{先写 $a+kT/2$；若有一轴一中心，再试 $f(x+2(b-a))=2c-f(x)$}。不必借助高中以外的数学名词或技巧。
\end{trapbox}
